% Corrigé intégré automatiquement pour la banque Exercices.
% Source : A_Integrer/DDS_03/050_Geometrie/050_Geometrie.tex
\begin{corrigebox}[Corrigé]
\CorrigeQuestion{1}
On écrit la fermeture géométrique sous forme vectorielle, puis on projette dans une base adaptée. Les produits scalaires donnent les relations de longueur et les produits vectoriels les directions normales ; les équations obtenues sont ensuite résolues avec les conditions géométriques du mécanisme.
\CorrigeQuestion{2}
La résolution consiste à identifier les données et l'inconnue, choisir la loi du cours adaptée, écrire l'équation symbolique avant toute application numérique, puis vérifier l'homogénéité, le signe et l'ordre de grandeur du résultat. Les hypothèses utilisées doivent être explicitement contrôlées à la fin.
\end{corrigebox}
